\documentclass[../../../include/open-logic-section]{subfiles}

\begin{document}

\olfileid{sth}{cardinals}{cp}
\olsection{کانتور دا اصول}

\olref[ordinals][vn]{sec} چیتے کرو۔ اوتھے اسیں حُسنِ ترتیب والے سیٹاں دی گل کردے ساں تے آکھیا سی کہ چنگا ہووے جے ساڈے کول ایہو جیہیاں شےواں ہون جیہڑیاں حُسنِ ترتیب والیاں ساختاں دی نمائندگی کرن۔ ایس گل نوں سامنے رکھ کے اسیں ترتیبی عدد پیش کیتے، تے فیر \olref[ordinals][ordtype]{ordtypesworklikeyouwant} وچ وکھایا کہ اوہ اینج ای کم کردے نیں جیویں اسیں چاہندے ساں، یعنی:
\[
	\ordtype{A, <} = \ordtype{B, \lessdot}
	\text{ اُدوں تے صرف اُدوں جدوں } \tuple{A, <} \isomorphic \tuple{B, \lessdot}.
\]
ہُن ہور پچھے \olref[sfr][siz][equ]{sec} چیتے کرو۔ اوتھے اسیں سادہ طریقے نال کم کردیاں سیٹ دی «وڈائی» دا خیال پیش کیتا سی۔ خاص طور اُتے اسیں آکھیا سی کہ دو سیٹ ہم تعداد نیں، $\cardeq{A}{B}$، اُدوں تے صرف اُدوں جدوں اک !!a{bijection} $f \colon A \to
B$ اے۔ ایہ خیال اپنے آپ وچ حُسنِ ترتیب دے خیال نالوں {سوکھا} اے۔ ساڈا کم صرف !!{bijection}s نال اے۔ ایہناں !!{bijection}s نے (ترتیب دے ہم ساخت نقشیاں وانگ) «کوئی ساخت وی سنبھال کے رکھنی» اے یا نہیں، ایہ گل ایتھے ضروری نہیں۔

ایہ ساریاں گلاں توں اک سِدھا خیال نکلدا اے۔ جیویں اسیں حُسنِ ترتیب والیاں ساختاں دا ماپ لین لئی کجھ شےواں، یعنی \emph{ترتیبی عدد}، پیش کیتے سن، اویویں اسیں وڈائی دا ماپ لین لئی کجھ شےواں، یعنی \emph{کارڈینل عدد}، پیش کر سکدے آں۔ ایس باب دا مقصد ایہو اے۔

ایہ دسن توں پہلاں کہ ایہ کارڈینل عدد کیہ ہون گے، سانوں اک اصول دَسنا چاہیدا اے جیہدے اُتے اوہ پورے اُترن۔ $\card{X}$ نوں سیٹ $X$ دی کارڈینیلٹی لئی ورتدے ہوئیاں، اسیں چاہواں گے کہ:
\[
	\card{A} = \card{B} \text{ اُدوں تے صرف اُدوں جدوں } \cardeq{A}{B}.
\]
اسیں اینوں \emph{کانتور دا} اصول آکھاں گے، کیوں جے شاید کانتور پہلا سی جیہدے ذہن وچ ایہ گل بڑی صاف سی۔ (ایہدا \emph{ہیوم دے} اصول نال تعلق \olref[hp]{sec} وچ ہور سمجھاواں گے۔) سو ساڈا مقصد $\card{X}$ دی تعریف ہر $X$ لئی ایہو جیہی دینا اے جیہدے توں کانتور دا اصول ملے۔

\end{document}
